% Fuente única de clase: Recursos/2025-09-09.pdf, páginas 1--10.
En la clase 7 pasamos del teorema de Cauchy a los tres teoremas de Sylow. Las acciones sobre clases laterales permiten estudiar la contención y el número de subgrupos de Sylow. Aplicamos estos resultados para descartar grupos simples de ciertos órdenes y para probar que algunos grupos de orden $pq$ son cíclicos.

Los fragmentos originales de \citet{aluffi2009} reproducidos en las pizarras corresponden a III.2: afirmación III.2.2, definición III.2.3, teorema III.2.5, proposición III.2.6 y resultados III.2.8--16. Incluyen las pruebas y los ejemplos desarrollados aquí. La comprobación de puntos fijos, propuesta como ejercicio III.1.1 en la página 3, se presenta con solución antes de usarla en el segundo teorema de Sylow. Los localizadores se verificaron en los fragmentos originales de las pizarras.

%%%%%%%%%%%%%%%%%%%%%%%%%%%%%%%%%%%%%%%%%%%%%%%%%%%%%%%%%%%%%%%%%%%%%%%%%%%%%%%
\section{De Cauchy a la existencia de subgrupos de Sylow}
%%%%%%%%%%%%%%%%%%%%%%%%%%%%%%%%%%%%%%%%%%%%%%%%%%%%%%%%%%%%%%%%%%%%%%%%%%%%%%%
%%%%%%%%%%%%%%%%%%%%%%%%%%%%%%%%%%%%%%%%%%%%%%%%%%%%%%%%%%%%%%%%%%%%%%%%%%%%%%%
\subsection{El número de subgrupos de orden primo}
%%%%%%%%%%%%%%%%%%%%%%%%%%%%%%%%%%%%%%%%%%%%%%%%%%%%%%%%%%%%%%%%%%%%%%%%%%%%%%%
% P01.
\begin{prop}\label{prop:C7-orden-primo}
Si el primo $p$ divide a $|G|$, con $G$ finito, el número $N$ de subgrupos cíclicos de orden $p$ satisface $N\equiv1\pmod p$.
\end{prop}
\begin{proof}
Retomemos el conjunto de tuplas de la prueba de Cauchy y sus puntos fijos bajo rotación:
\[
S=\set{(a_1,\ldots,a_p)\in G^p\given a_1\cdots a_p=e_G},\qquad
Z=\set{(a,\ldots,a)\given a^p=e_G}.
\]
Por el \cref{coro:C6-fijos}, $|Z|\equiv|S|\equiv0\pmod p$. Además de la tupla neutra, cada elemento de $Z$ corresponde a un elemento de orden exactamente $p$. Un grupo cíclico de orden $p$ tiene $p-1$ elementos no neutros, todos generadores. Dos subgrupos de orden $p$ distintos sólo comparten el neutro: un elemento común no neutro generaría ambos. Por ello,
\[
|Z|=1+N(p-1).
\]
Al reducir módulo $p$ obtenemos $0\equiv1-N$, y por tanto $N\equiv1\pmod p$.
\end{proof}
El recuento también puede expresarse con una acción por potencias. Sea $U_p$ el grupo multiplicativo de clases no nulas módulo $p$. Actúa en $Z$ mediante
\[
[m]\cdot(a,\ldots,a)=(a^m,\ldots,a^m).
\]
La regla no depende del representante porque $a^p=e_G$, y la composición de las potencias corresponde al producto de las clases. Cada $[m]\ne[0]$ tiene un inverso módulo $p$, de modo que actúa biyectivamente. La órbita de una tupla no neutra reúne los $p-1$ generadores de $\generado{a}$, pues $a^m=a$ implica $m\equiv1\pmod p$. Así se recupera la misma igualdad contando órbitas.
\begin{nota}[Corrección de la acción por potencias]
En la página 1 se escribe $\mathbb Z/(p-1)\mathbb Z$ actuando por $[m]\cdot(a,\ldots,a)=(a^m,\ldots,a^m)$. Con la suma como operación y esa fórmula, no es una acción: $[0]$ enviaría toda tupla a la neutra. La fórmula anterior utiliza $U_p$, con multiplicación, cuyo neutro es $[1]$. No se necesita identificarlo con un grupo cíclico para este recuento. Para $p>2$, su único punto fijo es la tupla neutra; para $p=2$, la acción es trivial y todas las tuplas son fijas, aunque el recuento $1+N(p-1)$ sigue siendo válido.
\end{nota}

%%%%%%%%%%%%%%%%%%%%%%%%%%%%%%%%%%%%%%%%%%%%%%%%%%%%%%%%%%%%%%%%%%%%%%%%%%%%%%%
\subsection{Primer teorema de Sylow}
%%%%%%%%%%%%%%%%%%%%%%%%%%%%%%%%%%%%%%%%%%%%%%%%%%%%%%%%%%%%%%%%%%%%%%%%%%%%%%%
% P02--P03.
\begin{defi}
Escribamos $|G|=p^r m$, con $p$ primo, $r\ge0$ y $p\nmid m$. Un \emph{$p$-subgrupo de Sylow} de $G$ es un subgrupo de orden $p^r$. Equivalentemente, es un $p$-subgrupo $P\le G$ cuyo índice $[G:P]$ no es divisible por $p$.
\end{defi}
Si $p\nmid|G|$, tenemos $r=0$ y el grupo trivial es un $p$-subgrupo de Sylow. Para demostrar la existencia en todos los casos, probamos primero una afirmación más general.

\begin{prop}\label{prop:C7-potencias}
Si $p^k\mid|G|$, donde $G$ es finito, $p$ es primo y $k\ge0$, entonces $G$ contiene un subgrupo de orden $p^k$.
\end{prop}
\begin{proof}
Hacemos inducción fuerte sobre $|G|$, para todos los exponentes admisibles. Si $k=0$, sirve $\{e_G\}$. Si $k=1$, el \cref{teo:C6-cauchy} proporciona un elemento de orden $p$, que genera el subgrupo buscado. Podemos suponer $k>1$.

Si existe un subgrupo propio $H<G$ con $p\nmid[G:H]$, la igualdad $|G|=[G:H]|H|$ implica $p^k\mid|H|$. Como $|H|<|G|$, la hipótesis inductiva aplicada a $H$ da un subgrupo de orden $p^k$.

Supongamos ahora que $p$ divide al índice de todo subgrupo propio de $G$. Para cada $a\notin Z(G)$, el centralizador $Z_G(a)$ es propio, así que su índice es divisible por $p$. La ecuación de clases da
\[
|G|=|Z(G)|+\sum_{a\in R}[G:Z_G(a)]
\quad\Longrightarrow\quad p\mid|Z(G)|.
\]
Por Cauchy existe $a\in Z(G)$ de orden $p$. El subgrupo $N=\generado{a}$ es central y, por tanto, normal: cada conjugación fija sus elementos. Además $|N|=p$, y
\[
|G/N|=\frac{|G|}{p}<|G|,\qquad p^{k-1}\mid|G/N|.
\]
La inducción produce un subgrupo $\bar P\le G/N$ de orden $p^{k-1}$. Sea $\pi:G\to G/N$ la proyección y tomemos $P=\pi^{-1}(\bar P)$:
\[
\begin{tikzcd}[column sep=large,row sep=large]
P=\pi^{-1}(\bar P)\arrow[r,two heads,destacado,"\restr{\pi}{P}"]\arrow[d,hook]&\bar P\arrow[d,hook]\\
G\arrow[r,two heads,"\pi"']&G/N.
\end{tikzcd}
\]
El conjunto $P$ es subgrupo por el \cref{lema:C4-preimagen}; contiene a $N$ y la restricción de $\pi$ es sobreyectiva con núcleo $N$. En consecuencia $P/N\cong\bar P$ y
\[
|P|=|N|\,|P/N|=p\,p^{k-1}=p^k.\qedhere
\]
\end{proof}
\begin{nota}[La preimagen en el paso inductivo]
En la página 3, después de obtener un subgrupo del cociente, aparece $G/N=P/N$. En general se necesita $\bar P=P/N\le G/N$, con $P=\pi^{-1}(\bar P)$; la igualdad con todo $G/N$ sólo corresponde al caso $\bar P=G/N$. El diagrama hace explícita la preimagen que permite levantar el subgrupo y calcular su orden.
\end{nota}

\begin{teo}[Primer teorema de Sylow]\label{teo:C7-sylow-I}
Todo grupo finito contiene un $p$-subgrupo de Sylow para cada primo $p$.
\end{teo}
\begin{proof}
Escribimos $|G|=p^r m$ con $p\nmid m$. La \cref{prop:C7-potencias} aplicada a $k=r$ da un subgrupo de orden $p^r$, que es de Sylow por definición.
\end{proof}

%%%%%%%%%%%%%%%%%%%%%%%%%%%%%%%%%%%%%%%%%%%%%%%%%%%%%%%%%%%%%%%%%%%%%%%%%%%%%%%
\section{Contención, normalizadores y número de subgrupos}
%%%%%%%%%%%%%%%%%%%%%%%%%%%%%%%%%%%%%%%%%%%%%%%%%%%%%%%%%%%%%%%%%%%%%%%%%%%%%%%
%%%%%%%%%%%%%%%%%%%%%%%%%%%%%%%%%%%%%%%%%%%%%%%%%%%%%%%%%%%%%%%%%%%%%%%%%%%%%%%
\subsection{Puntos fijos y segundo teorema de Sylow}
%%%%%%%%%%%%%%%%%%%%%%%%%%%%%%%%%%%%%%%%%%%%%%%%%%%%%%%%%%%%%%%%%%%%%%%%%%%%%%%
% P03--P04.
\begin{ejer}[Puntos fijos de un $p$-grupo]\label{ejer:C7-fijos}
Sea $H$ un $p$-grupo que actúa en un conjunto finito $X$ de cardinal no divisible por $p$. Demostrar que la acción tiene un punto fijo.
\end{ejer}
\begin{sol}
El \cref{coro:C6-fijos} da $|X^H|\equiv|X|\pmod p$. Como $|X|\not\equiv0\pmod p$, también $|X^H|\not\equiv0\pmod p$; en particular $|X^H|\ne0$. Existe entonces un elemento de $X$ fijado por todo $H$. En términos de la ecuación de clases, las órbitas no triviales aportan múltiplos de $p$, de modo que no pueden reunir por sí solas todos los elementos de $X$.
\end{sol}
En el ejercicio III.1.1 reproducido en la página 3, la condición sobre $|X|$ se expresa mediante congruencias; aquí hacemos explícita la finitud del conjunto para aplicar el recuento.

\begin{teo}[Segundo teorema de Sylow]\label{teo:C7-sylow-II}
Si $P$ es un $p$-subgrupo de Sylow de un grupo finito $G$ y $H\le G$ es un $p$-subgrupo, existe $g\in G$ tal que $H\subseteq gPg^{-1}$.
\end{teo}
\begin{proof}
Hacemos actuar $H$ en el conjunto de clases izquierdas $G/P$ mediante
\[
h\cdot(gP)=(hg)P.
\]
El conjunto tiene $[G:P]$ elementos, y este número no es divisible por $p$. Por el \cref{ejer:C7-fijos}, existe una clase $gP$ fija bajo todo $H$. Así, para cada $h\in H$,
\[
hgP=gP\quad\Longleftrightarrow\quad
g^{-1}hgP=P\quad\Longleftrightarrow\quad g^{-1}hg\in P.
\]
Por tanto $g^{-1}Hg\subseteq P$, es decir, $H\subseteq gPg^{-1}$.
\end{proof}

\begin{coro}\label{coro:C7-conjugados}
Todos los $p$-subgrupos de Sylow de $G$ son conjugados. En particular, uno de ellos es normal si y sólo si es el único.
\end{coro}
\begin{proof}
Si $Q$ también es de Sylow, el teorema da $Q\subseteq gPg^{-1}$. Ambos tienen el mismo orden $p^r$, así que la inclusión es una igualdad. Si $P$ es normal, todos sus conjugados son $P$ y no hay otros subgrupos de Sylow. Recíprocamente, cada conjugado de $P$ tiene el mismo orden, luego es de Sylow; si éste es único, todos coinciden con $P$, que es entonces normal.
\end{proof}

%%%%%%%%%%%%%%%%%%%%%%%%%%%%%%%%%%%%%%%%%%%%%%%%%%%%%%%%%%%%%%%%%%%%%%%%%%%%%%%
\subsection{Crecimiento dentro del normalizador}
%%%%%%%%%%%%%%%%%%%%%%%%%%%%%%%%%%%%%%%%%%%%%%%%%%%%%%%%%%%%%%%%%%%%%%%%%%%%%%%
% P04--P06.
\begin{lema}\label{lema:C7-normalizador}
Si $H$ es un $p$-subgrupo de un grupo finito $G$, entonces
\[
[N_G(H):H]\equiv[G:H]\pmod p.
\]
\end{lema}
\begin{proof}
Si $H=\{e_G\}$, su normalizador es $G$ y los índices son iguales. En general, hacemos actuar $H$ por multiplicación izquierda en $X=G/H$. Una clase $gH$ es fija si y sólo si $hgH=gH$ para todo $h\in H$, lo que equivale a $g^{-1}Hg\subseteq H$. Como estos dos subgrupos tienen el mismo orden finito, la inclusión es igualdad. Por tanto,
\[
gH\in X^H\quad\Longleftrightarrow\quad g\in N_G(H),\qquad
X^H=\set{gH\given g\in N_G(H)}.
\]
Los puntos fijos son exactamente las clases izquierdas de $H$ dentro de su normalizador, de modo que $|X^H|=[N_G(H):H]$. La congruencia $|X^H|\equiv|X|\pmod p$ prueba la afirmación.
\end{proof}

\begin{prop}\label{prop:C7-crecimiento}
Si $H$ es un $p$-subgrupo de $G$ que no es de Sylow, existe un $p$-subgrupo $H'$ con $H\trianglelefteq H'$ y $[H':H]=p$.
\end{prop}
\begin{proof}
Como $H$ no es de Sylow, $p\mid[G:H]$. El lema anterior implica $p\mid[N_G(H):H]$. Puesto que $H\trianglelefteq N_G(H)$, podemos formar el grupo $N_G(H)/H$, cuyo orden es divisible por $p$. Cauchy da en él un subgrupo cíclico $C$ de orden $p$.

Tomemos la preimagen $H'$ de $C$ bajo la proyección $N_G(H)\to N_G(H)/H$. Entonces $H\le H'\le N_G(H)$ y $H'/H\cong C$. Como $H$ es normal en su normalizador, también lo es en $H'$. El índice y el orden satisfacen
\[
[H':H]=|C|=p,\qquad |H'|=p|H|.
\]
Así, si $|H|=p^k$, tenemos $|H'|=p^{k+1}$: también $H'$ es un $p$-subgrupo.
\end{proof}

Supongamos que ya tenemos una cadena de $p$-subgrupos $\{e_G\}=H_0\le H_1\le\cdots\le H_k$, con $|H_i|=p^i$. Si $|G|=p^r m$ y $p\nmid m$, el segundo teorema sitúa $H_k$ en un subgrupo de Sylow. La proposición anterior permite además prolongar la cadena paso a paso:
\[
\underbrace{\{e_G\}=H_0<H_1<\cdots<H_k}_{|H_i|=p^i}
<H_{k+1}<\cdots<H_r,\qquad |H_r|=p^r.
\]
Cada nuevo paso tiene índice $p$ y el subgrupo anterior es normal en el siguiente.

%%%%%%%%%%%%%%%%%%%%%%%%%%%%%%%%%%%%%%%%%%%%%%%%%%%%%%%%%%%%%%%%%%%%%%%%%%%%%%%
\subsection{Tercer teorema de Sylow}
%%%%%%%%%%%%%%%%%%%%%%%%%%%%%%%%%%%%%%%%%%%%%%%%%%%%%%%%%%%%%%%%%%%%%%%%%%%%%%%
% P07.
\begin{teo}[Tercer teorema de Sylow]\label{teo:C7-sylow-III}
Si $|G|=p^r m$ y $p\nmid m$, el número $n_p$ de $p$-subgrupos de Sylow satisface
\[
n_p\mid m,\qquad n_p\equiv1\pmod p.
\]
\end{teo}
\begin{proof}
Fijemos un $p$-subgrupo de Sylow $P$. Por el segundo teorema todos los demás son sus conjugados; por el \cref{lema:C6-conjugados},
\[
n_p=[G:N_G(P)].
\]
La multiplicatividad de índices da
\[
m=[G:P]=[G:N_G(P)]\,[N_G(P):P]
=n_p[N_G(P):P].
\]
De aquí $n_p\mid m$. Por el \cref{lema:C7-normalizador},
\[
[N_G(P):P]\equiv[G:P]=m\pmod p.
\]
Sustituyendo en la igualdad de índices obtenemos $m\equiv n_pm\pmod p$. Como $p\nmid m$, podemos cancelar $m$ módulo $p$, y resulta $n_p\equiv1\pmod p$.
\end{proof}
Usamos $n_p$ para el número de subgrupos y $N_G(P)$ para el normalizador; esta distinción evita confundir dos usos de la letra $N$ en los recuentos.

%%%%%%%%%%%%%%%%%%%%%%%%%%%%%%%%%%%%%%%%%%%%%%%%%%%%%%%%%%%%%%%%%%%%%%%%%%%%%%%
\section{Aplicaciones a simplicidad y grupos de orden \texorpdfstring{$pq$}{pq}}
%%%%%%%%%%%%%%%%%%%%%%%%%%%%%%%%%%%%%%%%%%%%%%%%%%%%%%%%%%%%%%%%%%%%%%%%%%%%%%%
%%%%%%%%%%%%%%%%%%%%%%%%%%%%%%%%%%%%%%%%%%%%%%%%%%%%%%%%%%%%%%%%%%%%%%%%%%%%%%%
\subsection{Subgrupos normales forzados por el orden}
%%%%%%%%%%%%%%%%%%%%%%%%%%%%%%%%%%%%%%%%%%%%%%%%%%%%%%%%%%%%%%%%%%%%%%%%%%%%%%%
% Definición en P02; aplicaciones P07--P09.
\begin{defi}
Un grupo es \emph{simple} si es no trivial y carece de subgrupos normales propios no triviales.
\end{defi}

\begin{prop}\label{prop:C7-no-simple}
Si $|G|=mp^r$, con $p$ primo, $r\ge1$ y $1<m<p$, entonces $G$ no es simple.
\end{prop}
\begin{proof}
Por el tercer teorema, $n_p\mid m$ y $n_p\equiv1\pmod p$. Como $1\le n_p\le m<p$, se sigue $n_p=1$. El único $p$-subgrupo de Sylow es normal; tiene orden $p^r>1$ y es propio porque su índice es $m>1$. Es, por tanto, un subgrupo normal no trivial y propio.
\end{proof}
\begin{nota}[Hipótesis sobre el exponente]
La afirmación III.2.12 reproducida en la página 7 escribe el orden como $mp^r$ con $1<m<p$. Para concluir mediante un subgrupo normal no trivial se requiere $r\ge1$, que dejamos explícito. Si se permitiera $r=0$, el grupo simple $C_2$ tendría orden $2\cdot3^0$ y cumpliría $1<2<3$, contradiciendo el enunciado sin esa precisión.
\end{nota}

\begin{ejem}[Orden $2002$]
Como $2002=2\cdot7\cdot11\cdot13$, el número $n_{11}$ divide a $182=2\cdot7\cdot13$ y es congruente con $1$ módulo $11$. Sus posibles divisores y sus residuos son
\[
\begin{array}{c|rrrrrrrr}
d&1&2&7&13&14&26&91&182\\\hline
d\bmod11&1&2&7&2&3&4&3&6
\end{array}
\]
Sólo $1$ cumple ambas condiciones. Así existe un único subgrupo de Sylow de orden $11$, que es normal, propio y no trivial. Ningún grupo de orden $2002$ es simple.
\end{ejem}

\begin{ejem}[Orden $12$]
Para $|G|=12=2^2\cdot3$, las condiciones de Sylow dan
\[
n_2\in\{1,3\},\qquad n_3\in\{1,4\}.
\]
Si $n_3=1$, el subgrupo de orden $3$ es normal. Supongamos $n_3=4$ y denotemos esos subgrupos por $P_1,\ldots,P_4$. Como son distintos y tienen orden primo, sus intersecciones por pares son $\{e_G\}$. Por tanto,
\[
\left|\bigcup_{i=1}^4(P_i\setminus\{e_G\})\right|=4(3-1)=8.
\]
Quedan exactamente tres elementos distintos del neutro fuera de esa unión. Un $2$-subgrupo de Sylow tiene orden $4$ y ninguno de sus elementos puede tener orden $3$, por Lagrange. Sus tres elementos no neutros deben ser precisamente los tres restantes. Así, todo $2$-subgrupo de Sylow coincide con ese mismo conjunto de cuatro elementos, luego $n_2=1$ y es normal. En ambos casos existe un subgrupo normal no trivial y propio; $G$ no es simple.
\end{ejem}
\begin{nota}[Precisión del recuento]
En la página 9 se atribuye orden $2$ o $4$ a los tres elementos restantes. Lo justificamos al mostrar que pertenecen a un $2$-subgrupo de Sylow, completando el mismo argumento por recuento.
\end{nota}

\begin{ejem}[Orden $24$]
Si $|G|=24=2^3\cdot3$, entonces $n_2\mid3$ y $n_2\equiv1\pmod2$, de modo que $n_2=1$ o $3$. Si es $1$, el único subgrupo de orden $8$ es normal. Si es $3$, la conjugación actúa en el conjunto $X=\{P_1,P_2,P_3\}$ de subgrupos de Sylow y define
\[
\varphi:G\to\Sym(X)\cong S_3.
\]
Por el segundo teorema esta acción es transitiva y, al tener $X$ tres elementos, no es trivial. Así $\ker\varphi\ne G$. Tampoco puede ser trivial: de serlo, $\varphi$ sería inyectivo por la \cref{prop:C4-mono}, pero $24>6=|S_3|$. Su núcleo es entonces un subgrupo normal propio y no trivial, y $G$ no es simple.
\end{ejem}
La página 9 concluye que el núcleo es normal. Hemos completado la comprobación de que es propio y no trivial, necesaria para descartar la simplicidad.

%%%%%%%%%%%%%%%%%%%%%%%%%%%%%%%%%%%%%%%%%%%%%%%%%%%%%%%%%%%%%%%%%%%%%%%%%%%%%%%
\subsection{Un criterio de ciclicidad para el orden \texorpdfstring{$pq$}{pq}}
%%%%%%%%%%%%%%%%%%%%%%%%%%%%%%%%%%%%%%%%%%%%%%%%%%%%%%%%%%%%%%%%%%%%%%%%%%%%%%%
% P10.
\begin{prop}\label{prop:C7-pq}
Sea $G$ un grupo de orden $pq$, con $p<q$ primos. Si $q\not\equiv1\pmod p$, entonces $G$ es cíclico.
\end{prop}
\begin{proof}
Por Sylow, $n_p\mid q$ y $n_p\equiv1\pmod p$. Los únicos candidatos son $1$ y $q$; la hipótesis excluye $q$. Existe, por tanto, un único subgrupo $H$ de orden $p$, y es normal.

La conjugación preserva $H$ y define un homomorfismo
\[
\gamma:G\to\Aut(H),\qquad \gamma(g)(h)=ghg^{-1}.
\]
Como $H$ es cíclico de orden $p$, su grupo de automorfismos tiene $p-1$ elementos. En efecto, si $a$ genera $H$, cada automorfismo queda determinado por $a\mapsto a^i$, con $1\le i<p$, y cada una de esas elecciones define un automorfismo.

El orden de $\im\gamma$ divide a $|G|=pq$, por el primer teorema de isomorfismo y Lagrange, y también divide a $|\Aut(H)|=p-1$. Como $p$ y $q$ son mayores que cualquier divisor primo de $p-1$, tenemos $\mcd(pq,p-1)=1$. Así, $\im\gamma$ es trivial: cada conjugación fija los elementos de $H$, y $H\subseteq Z(G)$.

El índice $[G:Z(G)]$ divide a $[G:H]=q$, por lo que es $1$ o $q$. Así, $G/Z(G)$ es trivial o de orden primo; en ambos casos es cíclico. La solución del \cref{ejer:C6-centro} implica que $G$ es abeliano.

Por Cauchy existen $a,b\in G$ de órdenes $p$ y $q$. Como conmutan, $(ab)^{pq}=e_G$. Si $(ab)^t=e_G$, entonces $a^t=b^{-t}$ pertenece a $\generado{a}\cap\generado{b}$. El orden de esta intersección divide tanto a $p$ como a $q$, luego es trivial. Por ello $a^t=b^t=e_G$, lo que implica $p\mid t$ y $q\mid t$, y por tanto $pq\mid t$. El elemento $ab$ tiene orden $pq$ y genera todo $G$.
\end{proof}
\begin{nota}[Centro y pasos completados]
En la página 10 se escribe $H=Z(G)$ después de probar que la conjugación es trivial en $H$. La conclusión inmediata es $H\subseteq Z(G)$; el centro puede ser mayor. El argumento de índices anterior trata ambas posibilidades y conserva la estrategia del cociente por el centro. También hemos explicitado el recuento de $\Aut(H)$, por qué el núcleo del homomorfismo al grupo simétrico sirve en el ejemplo de orden $24$ y el paso final de abeliano a cíclico para el orden $pq$.
\end{nota}
